\subsection{Cohomological descent}
\label{subsec:dotted-handle-interpretation}

The dependence on a handle cocycle is multiplicative at the level of
character coupons.  We use this observation to prove the stronger statement
that adding a coboundary leaves every fixed-label graphical summand unchanged.
The full evaluator then descends immediately from \(Z^2(\mathcal H)\) to
\(H^2(\mathcal H)\).

\begin{proposition}[The dotted-handle operation]
\label{prop:dotted-handle-projector}
Let \(\delta\) be a dotted component of a one-ended Kirby presentation
\(\mathcal H\).  Choose a cooriented cross-section of the corresponding
\(1\)-handle, transverse to the framed attaching link.  For a fixed
simple-label state \(\lambda\), let \(B_\delta\) be the ordered tensor product
of the oriented labels meeting that cross-section, with the rigid dual used
at an oppositely oriented boundary occurrence.  The evaluator inserts
\[
  P_{\mathbf1,B_\delta}
  \in\operatorname{End}_{\cC}(B_\delta).
\]
If \(e_k\) is the oriented \(2\)-handle generator corresponding to
\(K_k\subset D_{\mathrm{fr}}\), then
\begin{equation}
  \partial_2(e_k)=\sum_\delta I_\delta(k)e_\delta,
  \label{eq:dotted-cellular-boundary}
\end{equation}
where \(I_\delta(k)\) is the signed intersection number of \(K_k\) with the
chosen cross-section of \(\delta\).
\end{proposition}

\begin{proof}
The first assertion is the dotted-circle rule in
\cref{def:dotted-circle-translation}, applied to the cut-open tangle of
\cref{def:cut-open-dotted-tangle}.  The second assertion is the standard
signed-intersection description of the cellular boundary map for a handle
decomposition.
\end{proof}

\begin{lemma}[Addition of character decorations]
\label{lem:addition-of-character-decorations}
Fix \(\mathcal H\), a simple-label state \(\lambda\), and the auxiliary data
\(\mathfrak a\).  For \(\omega,\omega'\in Z^2(\mathcal H)\), the decorated
graph for \(\omega+\omega'\) is obtained from the graph for \(\omega\) by
composing, on every framed component \(K_k\), the additional coupon
\((\eta^{\omega'_k})_{\lambda_k}\).  In particular, addition of cocycles acts
multiplicatively on the character decorations.
\end{lemma}

\begin{proof}
On \(K_k\), the required identity is
\[
  (\eta^{\omega_k+\omega'_k})_{\lambda_k}
  =
  (\eta^{\omega_k})_{\lambda_k}
  \circ
  (\eta^{\omega'_k})_{\lambda_k}.
\]
This holds because \(\eta_{\lambda_k}\) is an automorphism.  All other local
morphisms are unchanged.
\end{proof}

\begin{proposition}[Coboundaries act trivially on each fixed-label summand]
\label{prop:fixed-label-coboundary-triviality}
Let \(\omega\in Z^2(\mathcal H)\), let \(b\in C^1(\mathcal H)\), and fix a
simple-label state \(\lambda\).  Then
\[
  L^\eta_{\lambda,\omega+\delta^1b}(\mathcal H)
  =
  L^\eta_{\lambda,\omega}(\mathcal H)
  \quad\text{in }\operatorname{End}_{\cC}(\mathbf1).
\]
\end{proposition}

\begin{proof}
The group \(C^1(\mathcal H)\) is freely generated by the dotted
\(1\)-handles.  By \cref{lem:addition-of-character-decorations}, it is enough
to consider \(b=m e_\delta^*\) for one dotted component \(\delta\) and one
integer \(m\).

By \eqref{eq:dotted-cellular-boundary}, the extra exponent on the framed
component \(K_k\) is \(mI_\delta(k)\).  Monoidality of \(\eta\), together
with \(\eta_{X^*}=\eta_X^{-1}\), therefore assembles all additional coupons
meeting the cross-section into the single endomorphism
\[
  (\eta^m)_{B_\delta}\in\operatorname{End}_{\cC}(B_\delta).
\]
Here the orientation of each framed component fixes its label before the
cross-section is chosen; the intersection signs only record the variance of
its occurrences in \(B_\delta\).

The dotted-handle operation is \(P_{\mathbf1,B_\delta}\).  Since this
projector retains only the tensor-unit channel and
\((\eta^m)_{\mathbf1}=\operatorname{id}_{\mathbf1}\),
\cref{lem:unit-projector-absorbs-character} gives
\[
  P_{\mathbf1,B_\delta}\circ(\eta^m)_{B_\delta}
  =P_{\mathbf1,B_\delta}.
\]
Thus the additional coboundary decoration disappears before the common
exterior tangle is reglued.  Repeating this argument over the finite support
of \(b\) proves the claim.
\end{proof}

\begin{proposition}[Cohomological descent]
\label{prop:representative-independence}
If \(\omega,\omega'\in Z^2(\mathcal H)\) represent the same class in
\(H^2(\mathcal H;\mathbb Z)\), then
\[
  \CY_{\cC,\eta}(\mathcal H,\omega)
  =
  \CY_{\cC,\eta}(\mathcal H,\omega').
\]
\end{proposition}

\begin{proof}
Write \(\omega'-\omega=\delta^1b\).  By
\cref{prop:fixed-label-coboundary-triviality}, the two graphical morphisms
are equal for every label state \(\lambda\).  Their quantum-dimension weights
are identical, so summing over \(\lambda\) proves the equality.
\end{proof}

\begin{definition}[Evaluator on a cohomology class]
\label{def:cohomological-handle-evaluator}
For \(\xi\in H^2(\mathcal H;\mathbb Z)\), define
\[
  \CY_{\cC,\eta}(\mathcal H,\xi)
\]
to be \(\CY_{\cC,\eta}(\mathcal H,\omega)\) for any handle cocycle
\(\omega\) representing \(\xi\).  This is well defined by
\cref{prop:representative-independence}.  For a marked Kirby datum
\((\mathcal H,\beta)\), define
\[
  \CY_{\cC,\eta}(\mathcal H,\beta)
\]
by taking \(\xi=\xi_{\mathcal H,\beta}\) from
\cref{def:handle-cocycle}.
\end{definition}

\begin{remark}[Multiplicativity occurs before evaluation]
\label{rem:decoration-not-evaluator-multiplicativity}
The identity in \cref{lem:addition-of-character-decorations} concerns the
composition of coupons on one fixed labelled graph.  In general it does not
imply
\[
  \bigl\langle L^\eta_{\lambda,\omega+\omega'}(\mathcal H)\bigr\rangle
  =
  \bigl\langle L^\eta_{\lambda,\omega}(\mathcal H)\bigr\rangle
  \bigl\langle L^\eta_{\lambda,\omega'}(\mathcal H)\bigr\rangle
  \quad\text{or}\quad
  \CY_{\cC,\eta}(\mathcal H,\omega+\omega')
  =\CY_{\cC,\eta}(\mathcal H,\omega)
   \CY_{\cC,\eta}(\mathcal H,\omega').
\]
The conclusion proved here is instead that coboundary decorations act as the
identity after the dotted projectors are inserted, so the evaluator factors
through \(H^2(\mathcal H;\mathbb Z)\).
\end{remark}
